%%%%%%%%%%%%%%%%%%%%%%%%%%%%%%%%%%%%%%%%%%%%%%%%%%%%%%
\section{Experimental Pipeline}
%%%%%%%%%%%%%%%%%%%%%%%%%%%%%%%%%%%%%%%%%%%%%%%%%%%%%%

This chapter presents the experimental pipeline used to transform raw operational data from the HHLA Container Terminal Burchardkai into a structured event log, to discover process and simulation models from this data, and to evaluate the resulting models in a realistic operational setting. The objective is not to propose a new general framework for simulation trustworthiness or to solve the broader methodological question of simulation discovery in abstract terms. Rather, the aim is to instantiate an existing data-driven workflow in a real-world container-terminal context and to assess its applicability, strengths, and limitations for truck-handling operations. In this sense, the contribution of this thesis is empirical and case-oriented: it provides evidence on how process mining and simulation-discovery techniques can be applied in a complex logistics environment, rather than claiming to resolve the entire research field of simulation trustworthiness in general.

The experimental pipeline consists of several consecutive stages. First, raw operational data from different systems are integrated and cleaned. Second, a process-oriented event log is constructed by defining the case, deriving meaningful activities, selecting timestamps, and enriching the data with contextual information. Third, a process model is discovered from the event log using a process discovery algorithm. Fourth, this process model and the enriched event log are used as input for the discovery of a simulation model. Finally, the discovered simulation model is evaluated under baseline and scenario-based conditions. The following sections describe these stages in detail.

\subsection{Event Log Engineering}

The first stage of the experimental pipeline concerns the construction of a high-quality event log from heterogeneous operational data sources. In the present study, the raw data originate from the terminal operating environment and contain information about truck visits, container movements, yard states, and operational decisions. Since these data are stored in different systems and are not originally designed for process analysis, they require a dedicated engineering process before they can be used for process mining and simulation.

\subsubsection{Raw Data Sources}

The raw data used in this study come from two main operational perspectives. On the one hand, yard-related data describe the state and usage of the container yard, including storage locations, handling operations, and yard workload conditions. On the other hand, truck-related data capture the demand side of the process, including truck arrivals, appointment information, transport orders, and related operational events. These two data sources are highly relevant for the analysis of truck handling processes, but they are characterized by different temporal resolutions, naming conventions, and levels of detail. For this reason, a preprocessing step is required before the data can be integrated into a unified event log.

A key part of this step is the interpolation of missing or irregular observations. Yard interpolation is used to reconstruct incomplete yard states and to obtain a more continuous representation of yard conditions over time. Similarly, truck demand interpolation is applied to recover missing demand information and to ensure a regular temporal structure for subsequent analysis. These interpolation steps are necessary because process mining and simulation techniques require data that are sufficiently complete and temporally consistent to support the reconstruction of process behaviour.

\subsubsection{Data Integration and Cleaning}

To create a unified dataset, the raw data from the different sources are integrated using cross-system identifiers. This step is necessary because the original systems store information under different schemas, use different naming conventions, and record different subsets of operational events. In the integration phase, corresponding entities are identified, identifiers are harmonized, and records are aligned to a common temporal reference. The goal is to create a consistent representation of the same operational reality across different data sources.

After integration, the data are cleaned to remove inconsistencies and to improve their quality for downstream analysis. This includes the removal of duplicate records, the correction of formatting inconsistencies, the elimination of invalid or contradictory entries, and the filtering of observations that violate domain constraints. For example, records with impossible timestamps, illogical state transitions, or missing core information are removed or corrected where possible. The cleaning step is essential to ensure that the resulting event log reflects valid process executions and does not contain artifacts that could distort the discovered process model or the derived simulation behaviour.

\subsubsection{Case Definition}

Following data integration and cleaning, a process case must be defined. In this thesis, the process case is represented by a truck visit. A truck visit comprises the sequence of activities that occur from the moment the truck enters the terminal until it exits again. This case definition is appropriate because the focus of the study lies on truck-handling operations and their interaction with yard processes, resource availability, and terminal congestion. Defining the case at the level of the truck visit enables the reconstruction of a complete operational trace for each truck and allows the analysis of waiting times, turnaround times, and process variations.

\subsubsection{Activity Construction}

The raw operational records are then transformed into semantically meaningful process activities. Low-level system events are grouped into higher-level activities that correspond to operational steps recognizable from a process perspective. This abstraction step reduces the technical complexity of the data and improves the interpretability of the resulting event log. Examples include activities such as gate entry, yard arrival, container handling, and gate exit. In addition, resource information is aggregated where appropriate. Resource pooling is applied to reduce sparsity and to merge individual resources or resource classes into broader operational categories that are meaningful for process-level analysis. This step is particularly relevant in settings where the raw data contain a large number of highly specific resource identifiers that do not always produce robust process-level patterns.

\subsubsection{Timestamp Selection}

An important design choice in event log engineering is the selection of the correct timestamp for each activity. In many operational systems, multiple timestamps are available, such as the time at which a request was created, the time at which an activity was started, or the time at which it was completed. For process mining and simulation purposes, the most meaningful timestamp must be selected according to the intended interpretation of the process. In this study, the selected timestamp corresponds to the operational point in time that best reflects the actual execution of the activity within the truck-handling process. This decision is crucial because the ordering and timing of events directly affect the discovery of process behaviour, the calculation of waiting times, and the realism of the resulting simulation model.

\subsubsection{Waiting Time Generation}

To capture temporal inefficiencies and delays within the process, waiting times are derived from the reconstructed event sequences. Waiting time is defined as the delay between an activity becoming relevant and its actual execution. These waiting periods are of particular interest because they reflect operational congestion, resource contention, and queuing effects. By explicitly generating waiting-time information, the event log not only records the occurrence of activities but also captures the temporal dynamics that influence process performance. This is especially important for simulation purposes, where waiting times play a central role in reproducing realistic system behaviour.

\subsubsection{Context Enrichment}

To make the event log more informative and to support data-aware simulation, contextual attributes are added to the events. These attributes describe the operational context in which an activity occurs and allow the process to be analysed not only structurally but also with respect to its environmental conditions. The following categories of context information are incorporated:

\begin{itemize}
    \item \textbf{Container features}: information such as container type, size, priority, or destination.
    \item \textbf{Truck features}: information about the truck involved, including truck type, capacity, and service class.
    \item \textbf{Yard features}: attributes describing the storage area, such as yard zone, block occupancy, and available capacity.
    \item \textbf{Workload features}: indicators such as queue length, number of concurrent operations, or current system load.
    \item \textbf{Temporal features}: variables such as hour of the day, weekday, or shift.
\end{itemize}

These attributes enrich the event log and provide the basis for context-aware analysis. They are particularly valuable for simulation discovery because the behaviour of the process is often shaped by external conditions rather than by control-flow alone.

\subsubsection{Final Event Log}

The final event log is stored as a structured set of events containing, at minimum, a case identifier, an activity label, a timestamp, resource information, and relevant contextual attributes. The event log is exported in XES format to ensure compatibility with process mining tools and to enable reproducible analysis. In addition to the event-level representation, descriptive statistics are computed to characterize the observed process behaviour. These include the mean and maximum truck turnaround time, as well as distributions of waiting times, activity durations, and process frequencies. These statistics provide an initial overview of the data and support the subsequent process discovery and simulation stages.

\subsection{Process Discovery}

Once the event log has been constructed, the next stage of the pipeline concerns process discovery. The XES event log serves as input to a process discovery algorithm, in this study the Inductive Miner. The algorithm discovers a process model from the observed behaviour encoded in the event log and produces a process tree, which is subsequently transformed into a Petri net representation. The resulting Petri net reflects the control-flow structure of the observed truck-handling process and acts as the structural backbone for the subsequent simulation model.

The discovery configuration includes several parameters that influence the quality and complexity of the resulting model. These include the noise threshold, the treatment of infrequent behaviour, the selected event log input, and the structural settings of the discovery algorithm. The parameter choices are made in such a way that the resulting process model is both sufficiently expressive to capture the observed operational behaviour and simple enough to remain interpretable and usable for simulation purposes.

\subsection{Simulation Discovery}

The discovered Petri net and the enriched event log together form the input for the simulation-discovery stage. In this step, the process model is extended with the additional components required for executable simulation, including arrival behaviour, resource behaviour, timing information, and decision logic. The objective is to create a simulation model that is not only structurally consistent with the observed process but also capable of reproducing realistic operational performance.

\subsubsection{Parameter Discovery}

Simulation discovery requires the estimation of several parameter classes. These include control-flow parameters, arrival patterns, resource assignments, activity durations, waiting times, and decision rules. The control-flow component is derived from the discovered process model, while the temporal and behavioural components are estimated from the event log. In particular, the timing behaviour of activities is inferred from observed execution data, and the routing behaviour of the process is represented through data-aware decision rules where appropriate. This allows the simulation model to reproduce not only the activity sequences observed in reality, but also the temporal dynamics and contextual dependencies that shape operational performance.

\subsubsection{Discovery Configurations}

A number of configuration choices are made in the simulation-discovery stage. These include the selected sample size, the maximum tree depth used in the discovered decision structures, the specification of hyperparameters, the chosen training configurations, and the runtime settings used during model generation. These settings are not the primary contribution of the thesis, but they are essential for ensuring that the simulation model is both computationally feasible and sufficiently expressive. The selected configuration therefore reflects a balance between accuracy, complexity, and reproducibility.

\subsection{Simulation Experiments}

After the simulation model has been discovered, it is evaluated through a series of experiments designed to assess its usefulness in realistic operational scenarios. Two main experimental settings are considered.

\subsubsection{Baseline Replay}

In the baseline replay experiment, the simulation model is executed under conditions that reflect the observed historical process. The simulation outputs are compared with the original event log to assess whether the model is able to reproduce the control flow, timing behaviour, and performance patterns that were observed in reality. This experiment provides a first indication of whether the discovered model captures the essential structure of the process.

\subsubsection{Block Closure Scenario}

In addition to baseline replay, a scenario-based experiment is conducted to evaluate the behaviour of the simulation model under disruptive operating conditions. In particular, a block closure scenario is considered in which a yard block becomes temporarily unavailable or substantially constrained. This setting is highly relevant for container terminals because the temporary unavailability of a storage area can strongly affect truck routing, waiting times, yard workload, and overall terminal performance. The block closure scenario allows the simulation model to be tested under conditions that go beyond normal operations and to assess whether it can still produce meaningful and realistic results under altered constraints.

\subsection{Evaluation Logic}

The final stage of the pipeline concerns the evaluation of the discovered simulation model. The assessment is based on multiple dimensions rather than on a single measure. First, control-flow similarity is evaluated to determine whether the simulation model reproduces the observed process variants and activity sequences. Second, temporal similarity is assessed by comparing distributions of waiting times, activity durations, and overall turnaround times between the observed and simulated event logs. Third, congestion and arrival behaviour are evaluated to examine whether the model reproduces realistic workload dynamics. These evaluation measures are important because the usefulness of a simulation model does not depend only on structural correctness, but also on its ability to reproduce realistic operational behaviour under changing conditions.

The experimental pipeline therefore provides a coherent workflow from raw data to process discovery, simulation discovery, and evaluation. It serves as the methodological backbone of the empirical study and enables the investigation of how data-driven simulation models perform in a real container-terminal environment.
